\documentclass[10pt,a4paper]{amsart}
\usepackage{amsmath,amssymb,amsthm,graphicx}
\usepackage{verbatim}
\newenvironment{reference}{\comment}{\endcomment}
%\newenvironment{reference}{}{}
\newenvironment{slogan}{\comment}{\endcomment}
\newenvironment{history}{\comment}{\endcomment}
% For commutative diagrams we use Xy-pic
\usepackage[all]{xy}

% We use 2cell for 2-commutative diagrams.
\xyoption{2cell}
% Theorem environments.
%
\theoremstyle{plain}
\newtheorem{theorem}[subsection]{Theorem}
\newtheorem{proposition}[subsection]{Proposition}
\newtheorem{lemma}[subsection]{Lemma}

\theoremstyle{definition}
\newtheorem{definition}[subsection]{Definition}
\newtheorem{example}[subsection]{Example}
\newtheorem{exercise}[subsection]{Exercise}
\newtheorem{situation}[subsection]{Situation}

\theoremstyle{remark}
\newtheorem{remark}[subsection]{Remark}
\newtheorem{remarks}[subsection]{Remarks}

\numberwithin{equation}{subsection}

% Macros
%
\def\lim{\mathop{\mathrm{lim}}\nolimits}
\def\colim{\mathop{\mathrm{colim}}\nolimits}
\def\Spec{\mathop{\mathrm{Spec}}}
\def\Hom{\mathop{\mathrm{Hom}}\nolimits}
\def\Ext{\mathop{\mathrm{Ext}}\nolimits}
\def\SheafHom{\mathop{\mathcal{H}\!\mathit{om}}\nolimits}
\def\SheafExt{\mathop{\mathcal{E}\!\mathit{xt}}\nolimits}
\def\Sch{\mathit{Sch}}
\def\Mor{\mathop{\mathrm{Mor}}\nolimits}
\def\Ob{\mathop{\mathrm{Ob}}\nolimits}
\def\Sh{\mathop{\mathit{Sh}}\nolimits}
\def\NL{\mathop{N\!L}\nolimits}
\def\CH{\mathop{\mathrm{CH}}\nolimits}
\def\proetale{{pro\text{-}\acute{e}tale}}
\def\etale{{\acute{e}tale}}
\def\QCoh{\mathit{QCoh}}
\def\Ker{\mathop{\mathrm{Ker}}}
\def\Im{\mathop{\mathrm{Im}}}
\def\Coker{\mathop{\mathrm{Coker}}}
\def\Coim{\mathop{\mathrm{Coim}}}

% Boxtimes
%
\DeclareMathSymbol{\boxtimes}{\mathbin}{AMSa}{"02}

%
% Macros for moduli stacks/spaces
%
\def\QCohstack{\mathcal{QC}\!\mathit{oh}}
\def\Cohstack{\mathcal{C}\!\mathit{oh}}
\def\Spacesstack{\mathcal{S}\!\mathit{paces}}
\def\Quotfunctor{\mathrm{Quot}}
\def\Hilbfunctor{\mathrm{Hilb}}
\def\Curvesstack{\mathcal{C}\!\mathit{urves}}
\def\Polarizedstack{\mathcal{P}\!\mathit{olarized}}
\def\Complexesstack{\mathcal{C}\!\mathit{omplexes}}
% \Pic is the operator that assigns to X its picard group, usage \Pic(X)
% \Picardstack_{X/B} denotes the Picard stack of X over B
% \Picardfunctor_{X/B} denotes the Picard functor of X over B
\def\Pic{\mathop{\mathrm{Pic}}\nolimits}
\def\Picardstack{\mathcal{P}\!\mathit{ic}}
\def\Picardfunctor{\mathrm{Pic}}
\def\Deformationcategory{\mathcal{D}\!\mathit{ef}}
\usepackage{cleveref}
\crefname{originalreference}{original reference}{original references}
\Crefname{originalreference}{Original reference}{Original references}
\makeatletter
\newlabel{algebra-lemma-relations}{{algebra-lemma-relations}{source}{}{sourcepointer-1}{}}
\newlabel{algebra-lemma-relations@cref}{{[originalreference][0][]algebra-lemma-relations}{[1][0][]source}{}{}{}}
\makeatother
\title{A local factorisation criterion for a map followed by a surjection onto a flat module}
\author{The Stacks Project Authors}
\date{stacks-project:modules.tex:local-snapshot-20261012; GNU FDL 1.2 or later}
\begin{document}
\maketitle
\noindent\textit{Editorial note.} All delivered mathematical text is the author's own, copied from the pinned Stacks Project chapter source sources/stacks-modules/pinned-original-source.tex; no statement or proof step was rewritten, reordered or removed.\par\medskip
\noindent\textit{Editorial note.} Preamble provenance: the Stacks Project keeps one single shared preamble for every chapter; this card takes it directly from the staged sources/stacks-modules/preamble.tex as three verbatim spans (lines 8-12, 14-18 and 154-217), so no derived file was written and nothing was added or redefined. The wrapper class and the base packages are supplied by the assembler.\par\medskip
\noindent\textit{Editorial note.} Bibliography provenance: neither the counted statement nor its proof contains a citation (this lane rescans the whole span before writing the card), so the card carries no bibliography and no entry was reconstructed.\par\medskip
\noindent\textit{Editorial note.} Selection provenance: this card was opened under the file-based control contract (CONTROL.md), with the mechanical pre-card gate run on its statement and proof interval before the card was written.\par\medskip
\noindent\textit{Allowed standard prerequisites (editorial).}\begin{itemize}
\item The sheaf of ideals generated by finitely many local functions, and the interpretation of the sum of the tensor products of those functions with the given sections as a section of the tensor product with the ideal sheaf. Used at the start of the counted proof (source lines 2993-2996). Turns the hypothesis that the displayed complex is a complex into a relation in a tensor product.
\item The injectivity of the maps obtained by tensoring with a flat module, and the description of the tensor product of the ideal sheaf with a flat module by generators and relations. Used inside the counted proof (source lines 2996-3004), the last step by a cross-chapter pointer. Lifts the vanishing of the section of the tensor product to an explicit relation between the local functions and the local sections.
\end{itemize}\medskip
\ifdefined\hypertarget\hypertarget{sourcepointer-1}{}\fi
\noindent\textit{Source reference.} Original reference algebra-lemma-relations: sources/stacks-modules/pinned-original-source.tex, source lines 3004--3004. Navigation pointer only: the counted text refers to the Stacks item algebra-lemma-relations, which belongs to another chapter of the project and is not part of this card. The pointer carries the target's own label as its number because this chapter's staged snapshot does not contain the other chapter's tag table; its content is not reproduced here.\par\medskip
\begin{lemma}
\label{lemma-flat-eq}
Let $(X, \mathcal{O}_X)$ be a ringed space. Let $\mathcal{F}$ be a flat
$\mathcal{O}_X$-module. Let $U \subset X$ be open and let
$$
\mathcal{O}_U \xrightarrow{(f_1, \ldots, f_n)}
\mathcal{O}_U^{\oplus n} \xrightarrow{(s_1, \ldots, s_n)}
\mathcal{F}|_U
$$
be a complex of $\mathcal{O}_U$-modules. For every $x \in U$ there
exists an open neighbourhood $V \subset U$ of $x$ and a factorization
$$
\mathcal{O}_V^{\oplus n}
\xrightarrow{A}
\mathcal{O}_V^{\oplus m} \xrightarrow{(t_1, \ldots, t_m)}
\mathcal{F}|_V
$$
of $(s_1, \ldots, s_n)|_V$ such that $A \circ (f_1, \ldots, f_n)|_V = 0$.
\end{lemma}

\begin{proof}
Let $\mathcal{I} \subset \mathcal{O}_U$ be the sheaf of ideals
generated by $f_1, \ldots, f_n$. Then $\sum f_i \otimes s_i$ is
a section of $\mathcal{I} \otimes_{\mathcal{O}_U} \mathcal{F}|_U$
which maps to zero in $\mathcal{F}|_U$. As $\mathcal{F}|_U$ is flat
the map
$\mathcal{I} \otimes_{\mathcal{O}_U} \mathcal{F}|_U \to \mathcal{F}|_U$
is injective. Since $\mathcal{I} \otimes_{\mathcal{O}_U} \mathcal{F}|_U$
is the sheaf associated to the presheaf tensor product, we see
there exists an open neighbourhood $V \subset U$ of $x$ such
that $\sum f_i|_V \otimes s_i|_V$ is zero in
$\mathcal{I}(V) \otimes_{\mathcal{O}(V)} \mathcal{F}(V)$.
Unwinding the definitions using Algebra, Lemma \ref{algebra-lemma-relations}
we find $t_1, \ldots, t_m \in \mathcal{F}(V)$ and $a_{ij} \in \mathcal{O}(V)$
such that $\sum a_{ij}f_i|_V = 0$ and $s_i|_V = \sum a_{ij}t_j$.
\end{proof}

\end{document}
